\documentclass[10pt]{article}
\usepackage[utf8]{inputenc}
\usepackage[vietnamese]{babel}
\usepackage[letterpaper, top=1.6cm, bottom=1.6cm, left=1.8cm, right=1.8cm]{geometry}
\usepackage{amsmath, amssymb}
\usepackage{xcolor}
\usepackage{tikz}
\usepackage{array}
\usepackage{colortbl}
\usepackage{enumitem}
\usepackage{fancyhdr}

% Màu sắc chuẩn nhận diện Stewart Calculus
\definecolor{stewartcyan}{RGB}{0, 130, 195}
\definecolor{tablehead}{RGB}{236, 245, 252}
\definecolor{tableborder}{RGB}{140, 170, 195}
\definecolor{darktext}{RGB}{30, 30, 30}

\pagestyle{fancy}
\fancyhf{}
\fancyhead[L]{\textsf{\textbf{\large 82}}\qquad\textsf{\textbf{CHƯƠNG 2}}\quad\textsf{Giới hạn và Đạo hàm}}
\fancyhead[R]{}
\renewcommand{\headrulewidth}{0pt}
\fancyfoot[C]{%
    \fontsize{6.5pt}{8pt}\selectfont\color{black!70}%
    Bản quyền 2021 Cengage Learning. Bảo lưu mọi quyền. Không được sao chép, quét hoặc nhân bản toàn bộ hay từng phần. WCN 02-200-203\\
    Đánh giá biên tập cho thấy nội dung bị loại bỏ không ảnh hưởng nghiêm trọng đến trải nghiệm học tập tổng thể. Cengage Learning có quyền gỡ bỏ thêm nội dung bất kỳ lúc nào nếu yêu cầu hạn chế quyền sau đó đòi hỏi.
}
\renewcommand{\footrulewidth}{0pt}

% Biểu tượng máy tính vẽ đồ thị
\newcommand{\graphicon}{%
    \begin{tikzpicture}[scale=0.34, baseline=-0.3ex]
        \draw[stewartcyan, thick, rounded corners=0.8pt] (0,0) rectangle (0.85, 0.7);
        \draw[stewartcyan!60, thin] (0.08,0.35) -- (0.77,0.35);
        \draw[stewartcyan!60, thin] (0.42,0.08) -- (0.42,0.62);
        \draw[stewartcyan, thick] plot[smooth, tension=0.7] coordinates {(0.12,0.18) (0.3,0.58) (0.55,0.14) (0.73,0.52)};
    \end{tikzpicture}%
}

\setlist{nosep}

\begin{document}
\color{darktext}

% BANNER TIÊU ĐỀ 2.1 BÀI TẬP
\vspace*{-10pt}
\noindent
\begin{tikzpicture}[baseline=(current bounding box.north)]
    \useasboundingbox (0,0) rectangle (\linewidth, 0.85);
    \draw[stewartcyan, line width=1.1pt] (0, 0.82) -- (\linewidth, 0.82);
    \draw[stewartcyan, line width=0.5pt] (0, 0.74) -- (\linewidth, 0.74);
    
    \node[anchor=west, inner sep=0pt] at (0.05, 0.35) {\fontsize{22pt}{22pt}\selectfont\textbf{\textsf{\color{stewartcyan}2.1}}};
    \draw[black!45, line width=1.1pt] (1.35, 0.12) -- (1.35, 0.62);
    \node[anchor=west, inner sep=0pt] at (1.55, 0.35) {\fontsize{16pt}{16pt}\selectfont\textbf{\textsf{Bài tập}}};
    
    \draw[stewartcyan, line width=0.5pt] (3.65, 0.42) -- (\linewidth, 0.42);
    \draw[stewartcyan, line width=1.1pt] (3.65, 0.32) -- (\linewidth, 0.32);
\end{tikzpicture}

\vspace{4pt}

\noindent
\begin{minipage}[t]{0.485\textwidth}
\fontsize{8.7pt}{11pt}\selectfont

\textbf{1.} Một bồn chứa $1000\text{ gallon}$ nước, nước thoát ra từ đáy bồn trong nửa giờ. Các giá trị trong bảng dưới đây cho biết thể tích $V$ của lượng nước còn lại trong bồn (tính bằng gallon) sau $t$ phút.

\begin{center}
\vspace{-2pt}
\arrayrulecolor{tableborder}
\renewcommand{\arraystretch}{1.2}
\setlength{\tabcolsep}{6pt}
\begin{tabular}{|>{\columncolor{tablehead}}c|c|c|c|c|c|c|}
\hline
$t\text{ (phút)}$ & 5 & 10 & 15 & 20 & 25 & 30 \\ \hline
$V\text{ (gal)}$ & 694 & 444 & 250 & 111 & 28 & 0 \\ \hline
\end{tabular}
\vspace{-2pt}
\end{center}

\begin{enumerate}[label=\textbf{(\alph*)}, leftmargin=*]
    \item Nếu $P$ là điểm $(15, 250)$ trên đồ thị của $V$, hãy tìm hệ số góc của các cát tuyến $PQ$ khi $Q$ là điểm trên đồ thị với $t = 5, 10, 20, 25$ và $30$.
    \item Ước lượng hệ số góc của tiếp tuyến tại $P$ bằng cách lấy trung bình cộng hệ số góc của hai cát tuyến.
    \item Sử dụng đồ thị của $V$ để ước lượng hệ số góc của tiếp tuyến tại $P$. (Hệ số góc này biểu thị tốc độ nước chảy ra khỏi bồn sau 15 phút.)
\end{enumerate}

\vspace{5pt}
\textbf{2.} Một sinh viên mua một chiếc đồng hồ thông minh theo dõi số bước chân cô đi bộ trong ngày. Bảng dưới đây cho biết số bước chân ghi nhận được sau $t$ phút kể từ 3:00 PM vào ngày đầu tiên cô đeo đồng hồ.

\begin{center}
\vspace{-2pt}
\arrayrulecolor{tableborder}
\renewcommand{\arraystretch}{1.2}
\setlength{\tabcolsep}{8pt}
\begin{tabular}{|>{\columncolor{tablehead}}c|c|c|c|c|c|}
\hline
$t\text{ (phút)}$ & 0 & 10 & 20 & 30 & 40 \\ \hline
Số bước & 3438 & 4559 & 5622 & 6536 & 7398 \\ \hline
\end{tabular}
\vspace{-2pt}
\end{center}

\begin{enumerate}[label=\textbf{(\alph*)}, leftmargin=*]
    \item Tìm hệ số góc của các cát tuyến tương ứng với các khoảng thời gian đã cho của $t$. Các hệ số góc này biểu thị điều gì?\\
    \textbf{(i)} $[0, 40]$ \qquad\quad \textbf{(ii)} $[10, 20]$ \qquad\quad \textbf{(iii)} $[20, 30]$
    \item Ước lượng tốc độ đi bộ của sinh viên, tính theo số bước mỗi phút, vào lúc 3:20 PM bằng cách lấy trung bình cộng hệ số góc của hai cát tuyến.
\end{enumerate}

\vspace{5pt}
\textbf{3.} Điểm $P(2, -1)$ nằm trên đường cong $y = 1/(1 - x)$.
\begin{enumerate}[label=\textbf{(\alph*)}, leftmargin=*]
    \item Nếu $Q$ là điểm $(x, 1/(1 - x))$, hãy tìm hệ số góc của cát tuyến $PQ$ (làm tròn chính xác đến sáu chữ số thập phân) đối với các giá trị sau của $x$:
    \begin{center}
    \vspace{-1pt}
    \begin{tabular}{@{}llll@{}}
    \textbf{(i)} $1.5$ & \textbf{(ii)} $1.9$ & \textbf{(iii)} $1.99$ & \textbf{(iv)} $1.999$ \\
    \textbf{(v)} $2.5$ & \textbf{(vi)} $2.1$ & \textbf{(vii)} $2.01$ & \textbf{(viii)} $2.001$
    \end{tabular}
    \vspace{-1pt}
    \end{center}
    \item Sử dụng các kết quả từ phần (a), hãy dự đoán giá trị hệ số góc của tiếp tuyến với đường cong tại điểm $P(2, -1)$.
    \item Sử dụng hệ số góc từ phần (b), hãy tìm phương trình của tiếp tuyến với đường cong tại $P(2, -1)$.
\end{enumerate}

\vspace{5pt}
\textbf{4.} Điểm $P(0.5, 0)$ nằm trên đường cong $y = \cos \pi x$.
\begin{enumerate}[label=\textbf{(\alph*)}, leftmargin=*]
    \item Nếu $Q$ là điểm $(x, \cos \pi x)$, hãy tìm hệ số góc của cát tuyến $PQ$ (làm tròn chính xác đến sáu chữ số thập phân) đối với các giá trị sau của $x$:
    \begin{center}
    \vspace{-1pt}
    \begin{tabular}{@{}lll r@{}}
    \textbf{(i)} $0$ & \textbf{(ii)} $0.4$ & \textbf{(iii)} $0.49$ & \\
    \textbf{(iv)} $0.499$ & \textbf{(v)} $1$ & \textbf{(vi)} $0.6$ & \graphicon \\
    \textbf{(vii)} $0.51$ & \textbf{(viii)} $0.501$ & & 
    \end{tabular}
    \vspace{-1pt}
    \end{center}
    \item Sử dụng các kết quả từ phần (a), hãy dự đoán giá trị hệ số góc của tiếp tuyến với đường cong tại $P(0.5, 0)$.
\end{enumerate}

\end{minipage}%
\hfill
\begin{minipage}[t]{0.485\textwidth}
\fontsize{8.7pt}{11pt}\selectfont

\begin{enumerate}[label=\textbf{(\alph*)}, leftmargin=*]
    \setcounter{enumi}{2}
    \item Sử dụng hệ số góc từ phần (b), hãy tìm phương trình của tiếp tuyến với đường cong tại $P(0.5, 0)$.
    \item Vẽ phác thảo đường cong, hai trong số các cát tuyến, và tiếp tuyến.
\end{enumerate}

\vspace{5pt}
\textbf{5.} Mặt cầu của một cây cầu được treo cao 275 feet phía trên một dòng sông. Nếu một viên sỏi rơi từ thành cầu xuống, độ cao (tính bằng feet) của viên sỏi phía trên mặt nước sau $t$ giây được cho bởi phương trình $y = 275 - 16t^2$.
\begin{enumerate}[label=\textbf{(\alph*)}, leftmargin=*]
    \item Tìm vận tốc trung bình của viên sỏi trong khoảng thời gian bắt đầu từ khi $t = 4$ và kéo dài trong\\
    \textbf{(i)} 0.1 giây \qquad \textbf{(ii)} 0.05 giây \qquad \textbf{(iii)} 0.01 giây
    \item Ước lượng vận tốc tức thời của viên sỏi sau 4 giây.
\end{enumerate}

\vspace{5pt}
\textbf{6.} Nếu một hòn đá được ném thẳng đứng lên trên trên bề mặt Sao Hỏa với vận tốc $10\text{ m/s}$, độ cao của nó tính bằng mét sau $t$ giây được cho bởi $y = 10t - 1.86t^2$.
\begin{enumerate}[label=\textbf{(\alph*)}, leftmargin=*]
    \item Tìm vận tốc trung bình trong các khoảng thời gian đã cho:
    \begin{center}
    \vspace{-1pt}
    \begin{tabular}{@{}lll@{}}
    \textbf{(i)} $[1, 2]$ & \textbf{(ii)} $[1, 1.5]$ & \textbf{(iii)} $[1, 1.1]$ \\
    \textbf{(iv)} $[1, 1.01]$ & \textbf{(v)} $[1, 1.001]$ &
    \end{tabular}
    \vspace{-1pt}
    \end{center}
    \item Ước lượng vận tốc tức thời khi $t = 1$.
\end{enumerate}

\vspace{5pt}
\textbf{7.} Bảng dưới đây cho biết vị trí của một người lái xe mô tô sau khi tăng tốc từ trạng thái nghỉ.

\begin{center}
\vspace{-2pt}
\arrayrulecolor{tableborder}
\renewcommand{\arraystretch}{1.2}
\setlength{\tabcolsep}{5pt}
\begin{tabular}{|>{\columncolor{tablehead}}c|c|c|c|c|c|c|c|}
\hline
$t\text{ (giây)}$ & 0 & 1 & 2 & 3 & 4 & 5 & 6 \\ \hline
$s\text{ (feet)}$ & 0 & 4.9 & 20.6 & 46.5 & 79.2 & 124.8 & 176.7 \\ \hline
\end{tabular}
\vspace{-2pt}
\end{center}

\begin{enumerate}[label=\textbf{(\alph*)}, leftmargin=*]
    \item Tìm vận tốc trung bình cho mỗi khoảng thời gian:\\
    \textbf{(i)} $[2, 4]$ \qquad \textbf{(ii)} $[3, 4]$ \qquad \textbf{(iii)} $[4, 5]$ \qquad \textbf{(iv)} $[4, 6]$
    \item Sử dụng đồ thị của $s$ như một hàm số của $t$ để ước lượng vận tốc tức thời khi $t = 3$.
\end{enumerate}

\vspace{5pt}
\textbf{8.} Độ dời (tính bằng centimét) của một chất điểm chuyển động qua lại dọc theo một đường thẳng được xác định bởi phương trình chuyển động $s = 2\sin \pi t + 3\cos \pi t$, trong đó $t$ được đo bằng giây.
\begin{enumerate}[label=\textbf{(\alph*)}, leftmargin=*]
    \item Tìm vận tốc trung bình trong mỗi khoảng thời gian:
    \begin{center}
    \vspace{-1pt}
    \begin{tabular}{@{}p{3.2cm}p{3.2cm}@{}}
    \textbf{(i)} $[1, 2]$ & \textbf{(ii)} $[1, 1.1]$ \\[2pt]
    \textbf{(iii)} $[1, 1.01]$ & \textbf{(iv)} $[1, 1.001]$
    \end{tabular}
    \vspace{-1pt}
    \end{center}
    \item Ước lượng vận tốc tức thời của chất điểm khi $t = 1$.
\end{enumerate}

\vspace{5pt}
\textbf{9.} Điểm $P(1, 0)$ nằm trên đường cong $y = \sin(10\pi/x)$.
\begin{enumerate}[label=\textbf{(\alph*)}, leftmargin=*]
    \item Nếu $Q$ là điểm $(x, \sin(10\pi/x))$, hãy tìm hệ số góc của cát tuyến $PQ$ (làm tròn chính xác đến bốn chữ số thập phân) với $x = 2;\; 1.5;\; 1.4;\; 1.3;\; 1.2;\; 1.1;\; 0.5;\; 0.6;\; 0.7;\; 0.8$ và $0.9$. Các hệ số góc dường như có tiến đến một giới hạn nào không?
    \item Sử dụng đồ thị của đường cong để giải thích tại sao các hệ số góc của các cát tuyến ở phần (a) không ở gần với hệ số góc của tiếp tuyến tại $P$.
    \item Bằng cách chọn các cát tuyến phù hợp, hãy ước lượng hệ số góc của tiếp tuyến tại điểm $P$.
\end{enumerate}

\end{minipage}

\end{document}